% ======================================================================
\section{The observable and its expansion}\label{sec:observable}
% ======================================================================

\purpose{We define what is measured, a dressed gluon of momentum $k$ in the final state, and turn this definition into amplitudes built from the coefficients of $\Omega$ and from Wilson lines. The leading order is derived step by step, because every later manipulation is a longer version of the same three steps (act with $\bar\Omega^\dagger\Omega$ on the vacuum, project on the final state, square). The $O(g^4)$ terms are then sorted into three Groups and fifteen rows, which are the units of the rest of the paper.}

\subsection{Single inclusive gluon spectrum}\label{sec:observable_def}

The projectile, before the collision, is described by its valence state $\ket{\Psi_{\rm v}}$ times the dressing of the soft modes, $\Omega\ket{0}$, where $\ket0$ is the vacuum of the soft modes. The target is described by the eikonal scattering matrix $S$. At high energy the interaction of a soft gluon and of the valence charges with the target is a colour rotation by the adjoint Wilson line $U(\v x)$ at fixed transverse position:
\begin{align}\label{eq:Srotation}
S^\dagger\rho^a(\v x)S&=U^{ab}(\v x)\,\rho^b(\v x)\,,\notag\\
S^\dagger a^a_i(k^+,\v x)S&=U^{ab}(\v x)\,a^b_i(k^+,\v x)\,,\qquad
S^\dagger a^{\dagger a}_i(k^+,\v x)S=U^{ab}(\v x)\,a^{\dagger b}_i(k^+,\v x)\,.
\end{align}
After the collision the state is $S\,\Omega\ket0$. A gluon observed in the final state is a dressed (physical) gluon; the number of dressed gluons is counted by the bare number operator in the basis rotated back by $\Omega^\dagger$~\cite{Kovner:2006wr}. The single inclusive spectrum is therefore
\begin{align}\label{eq:observable}
\frac{d^3N}{d^3k}=\frac{1}{(2\pi)^3}\int_{w,w'}e^{-i\v k\cdot(\v w-\v w')}\,
\Big\langle\bra0\,\Omega^\dagger S^\dagger\Omega\;a^{\dagger a}_i(k^+,\v w')\,a^a_i(k^+,\v w)\;\Omega^\dagger S\,\Omega\ket0\Big\rangle\,,
\end{align}
where $\langle\dots\rangle$ denotes the average over the valence state, i.e.\ over the projectile charge configurations, taken at the end. The factor $(2\pi)^{-3}$ and the phase follow from $a^{\dagger a}_i(k)a^a_i(k)=\int_{w,w'}e^{-i\v k\cdot(\v w-\v w')}a^{\dagger a}_i(k^+,\v w')a^a_i(k^+,\v w)$ and the normalisation~\eqref{eq:aadagger}. The amplitude stands at $\v w$ and the conjugate amplitude at $\v w'$; we use this convention, $e^{-i\v k\cdot(\v w-\v w')}$ with the conjugated factor at the primed points, throughout the paper.

Inserting $SS^\dagger=1$ and defining $\bar\Omega\equiv S^\dagger\Omega S$, which is $\Omega$ with $\rho\to U\rho$ and $a\to Ua$ according to Eq.~\eqref{eq:Srotation}, we obtain
\begin{align}
\frac{d^3N}{d^3k}=\frac{1}{(2\pi)^3}\int_{w,w'}e^{-i\v k\cdot(\v w-\v w')}\,\Big\langle\bra0\,\Omega^\dagger\bar\Omega\;U^{ab'}(\v w')a^{\dagger b'}_i(k^+,\v w')\,U^{ab}(\v w)a^b_i(k^+,\v w)\;\bar\Omega^\dagger\Omega\ket0\Big\rangle\,.
\end{align}
A bar on any coefficient of $\Omega$ denotes, here and below, the replacement $\rho^a(\v u)\to U^{ab}(\v u)\rho^b(\v u)$ of every charge in it. Between $a(k^+,\v w)$ and $\bar\Omega^\dagger\Omega\ket0$ we insert a complete set of soft states. The vacuum does not contribute, and states with three or more gluons contribute at order $g^6$. Hence
\begin{align}\label{eq:N_amplitudes}
\frac{d^3N}{d^3k}=\frac{1}{(2\pi)^3}\Big\langle\mathcal A^{a\dagger}_i(k)\,\mathcal A^a_i(k)+\int dq^+\!\int_{v}\mathcal A^{ab\dagger}_{ij}(k;q^+,\v v)\,\mathcal A^{ab}_{ij}(k;q^+,\v v)\Big\rangle\,,
\end{align}
with the one-gluon and two-gluon amplitudes
\begin{align}\label{eq:amplitudes_def}
\mathcal A^a_i(k)&=\sqrt{2\pi}\int_{w}e^{-i\v k\cdot\v w}\,U^{ab}(\v w)\,\bra{g^b_i(k^+,\v w)}\bar\Omega^\dagger\Omega\ket0\,,\notag\\
\mathcal A^{ab}_{ij}(k;q^+,\v v)&=\sqrt{2\pi}\int_{w}e^{-i\v k\cdot\v w}\,U^{ac}(\v w)\,\bra{g^b_j(q^+,\v v)\,g^c_i(k^+,\v w)}\bar\Omega^\dagger\Omega\ket0\,.
\end{align}
The colour index $b$ of the second gluon in $\mathcal A^{ab}_{ij}$ is that of the bare state; since the second gluon is integrated over, its rotation by $U(\v v)$ drops out of the sum. The operators $\bar\Omega^\dagger\Omega$ are evaluated with the expansion~\eqref{eq:OmegaX}: $\Omega$ acting on $\ket0$ involves only $\mf N$, $\mf A_2$ and $\mf B_2$; $\bar\Omega^\dagger$ then acts on states with up to two gluons. Note that $\bar\Omega^\dagger\Omega=1$ for $U=1$, so that every amplitude vanishes in the absence of the target.

From here on $\mf A\equiv\mf A^{(1)}_2$ and $\mf A^\dagger\equiv\mf A^{(1)\dagger}_2$; the unitarity relation $\mf A^{(1)\dagger}_1(p^+,\v x)=-\mf A(p^+,-\v x)$ is used to eliminate $\mf A^{(1)}_1$ wherever possible, and $\mf A^{(3)}\equiv\mf A^{(3)}_2$.

\paragraph{Why the number operator is sandwiched between $\Omega^\dagger$ and $\Omega$.} After the collision the soft gluons that were part of the coherent cloud of the projectile are still attached to it unless the target has disturbed them. A detector does not count the bare quanta of this cloud; it counts gluons that are asymptotic states, i.e.\ eigenstates of the interacting Hamiltonian, $\Omega\,a^\dagger\ket0$. The projection on these states is what $\Omega\,a^\dagger a\,\Omega^\dagger$ in Eq.~\eqref{eq:observable} does. Two consequences will be used repeatedly: without a target ($S=1$) nothing is produced, to all orders, because $\Omega^\dagger\Omega=1$; and the final-state dressing brings the coefficients $\mf B_3$ and $\mf C_1$, which describe what happens to a gluon that is already present, into the amplitudes.

\subsection{Leading order, step by step}\label{sec:LO}

\paragraph{Step 1: $\bar\Omega^\dagger\Omega$ on the vacuum at order $g$.} To first order in $g$ the expansion~\eqref{eq:OmegaX} is
\begin{align}
\Omega=1+\int\frac{dp^+}{2\pi}\int_x\Big[\mf A^{(1)b}_{1j}(p^+,\v x)\,a^b_j(p^+,\v x)+\mf A^b_j(p^+,-\v x)\,a^{\dagger b}_j(p^+,\v x)\Big]+O(g^2)\,,
\end{align}
where the unitarity relation $-\mf A^{(1)\dagger}_1(p^+,\v x)=\mf A(p^+,-\v x)$ was used. In $\bar\Omega^\dagger$ the charges and the gluon operators are rotated, and the roles of creation and annihilation parts are exchanged. Acting on the vacuum, only the creation parts survive:
\begin{align}\label{eq:LOstate}
\bar\Omega^\dagger\Omega\ket0\Big|_{O(g)}=\int\frac{dp^+}{2\pi}\int_x\Big[\mf A^c_j(p^+,-\v x)-\bar{\mf A}^b_j(p^+,-\v x)\,U^{bc}(\v x)\Big]a^{\dagger c}_j(p^+,\v x)\ket0\,.
\end{align}
The first term is the gluon created by $\Omega$ before the collision; the second is the gluon that $\bar\Omega^\dagger$ removes, with the amplitude of the rotated charges.

\paragraph{Step 2: projection on the final state.} With the normalisation~\eqref{eq:states}, $\bra{g^b_i(k^+,\v w)}a^{\dagger c}_j(p^+,\v x)\ket0=\sqrt{2\pi}\,\delta^{bc}\delta_{ij}\,\delta(k^+-p^+)\,\delta^{(2)}(\v w-\v x)$. Inserting Eq.~\eqref{eq:LOstate} into the definition~\eqref{eq:amplitudes_def} and using $U^{ab}(\v w)U^{b'b}(\v w)=\delta^{ab'}$ in the second term gives
\begin{align}\label{eq:ALO}
\mathcal A^{a}_{{\rm LO}\,i}(k)&=\int_{w}e^{-i\v k\cdot\v w}\Big[U^{ab}(\v w)\,\mf A^b_i(k^+,-\v w)-\bar{\mf A}^a_i(k^+,-\v w)\Big]
\notag\\&=\frac{ig}{\sqrt2\,\pi\sqrt{k^+}}\int_{w,x}e^{-i\v k\cdot\v w}\,\Kv^i(\v x-\v w)\big[U(\v w)-U(\v x)\big]^{ab}\rho^b(\v x)\,.
\end{align}
In the second line the Weizs\"acker--Williams amplitude~\eqref{eq:A1} was inserted. The first term is the emission of the gluon before the collision, followed by the scattering of the gluon, which is rotated by $U(\v w)$ at its own position; the second is the emission after the collision by the charge at $\v x$, which has been rotated by $U(\v x)$.

\paragraph{Step 3: the square.} Since $\rho^a(\v x)$ is Hermitian, the conjugate amplitude is the same expression with the charge to the left; multiplying and dividing by $(2\pi)^3$, Eq.~\eqref{eq:N_amplitudes}, the LO spectrum is
\begin{align}\label{eq:LO}
\frac{d^3N_{\rm LO}}{d^3k}=\ &\frac{1}{(2\pi)^3}\,\frac{g^2}{2\pi^2}\,\frac{1}{k^+}\int_{x,x',w,w'}e^{-i\v k\cdot(\v w-\v w')}\,
\Kv(\v x-\v w)\cdot\Kv(\v x'-\v w')
\notag\\&\times
\Big\langle\big[U(\v x')-U(\v w')\big]^{ab'}\rho^{b'}(\v x')\,\big[U(\v x)-U(\v w)\big]^{ab}\rho^{b}(\v x)\Big\rangle\,.
\end{align}
Note that $\frac{1}{(2\pi)^3}\frac{g^2}{2\pi^2}=\frac{\alpha_s}{4\pi^4}$. The conjugated factor, at the primed points, stands to the left: this is the origin of the phase convention $e^{-i\v k\cdot(\v w-\v w')}$ and of the order $\rho(\v x')\rho(\v x)$ of the charges.

\paragraph{Physical interpretation of the leading order.} Three features of Eq.~\eqref{eq:LO} recur in every NLO term. (i) \emph{Production requires that the target resolves the gluon from its source.} The amplitude contains the difference $U(\v w)-U(\v x)$ of the Wilson lines at the position of the gluon and of the emitting charge. If the target rotates both in the same way, the Weizs\"acker--Williams cloud is rotated together with its source and remains a cloud: nothing is radiated. All NLO terms will be written with such differences. (ii) \emph{The projectile enters through a correlator of its charges, the target through a correlator of Wilson lines,} and the two are tied together only by transverse coordinates. This is what allows the rapidity logarithms of the NLO correction to act separately on the two factors. (iii) \emph{The longitudinal dependence is $1/k^+$,} i.e.\ the spectrum is flat in the rapidity of the gluon. All NLO terms have the same overall $1/k^+$; the logarithms $\log(\vee/k^+)$ and $\log(k^+/\Lambda)$ are the first sign of a dependence on the rapidity distances to the projectile and to the target.

\paragraph{Colour-singlet projectile.} For a projectile that is a colour singlet, or more generally whose charge correlator is colour diagonal,
\begin{align}\label{eq:singlet}
\big\langle\rho^a(\v x)\rho^b(\v y)\big\rangle=\frac{\delta^{ab}}{N_c^2-1}\,n(\v x,\v y)\,,\qquad n(\v x,\v y)=\big\langle\rho^f(\v x)\rho^f(\v y)\big\rangle\,,
\end{align}
and the LO spectrum becomes
\begin{align}\label{eq:LOsinglet}
\frac{d^3N_{\rm LO}}{d^3k}=\ &\frac{\alpha_s}{4\pi^4(N_c^2-1)}\,\frac{1}{k^+}\int_{x,y,w,w'}e^{-i\v k\cdot(\v w-\v w')}\,\Kv(\v x-\v w)\cdot\Kv(\v y-\v w')
\notag\\&\times\,n(\v x,\v y)\,
\mathcal O_{\rm LO}(\v x,\v y;\v w,\v w')\,,
\end{align}
with the target operator
\begin{align}\label{eq:OLO}
\mathcal O_{\rm LO}(\v x,\v y;\v w,\v w')={\rm Tr}\Big\{\big[U^\dagger(\v x)-U^\dagger(\v w)\big]\big[U(\v y)-U(\v w')\big]\Big\}\,.
\end{align}
Since $U$ is real, ${\rm Tr}\{A^\dagger B\}={\rm Tr}\{B^\dagger A\}$ for adjoint matrices, and for a symmetric $n$ the integrand is invariant under $(\v x,\v w)\leftrightarrow(\v y,\v w')$; in the colour-singlet form the sign of the exponent is therefore immaterial. Eq.~\eqref{eq:LOsinglet} is the formula of Kovchegov and Tuchin~\cite{Kovchegov:2001sc}: for a single quark, $n\to C_F$ times a $\delta$-function, it reproduces their normalisation $\alpha_sC_F/(4\pi^4)$ exactly. It is also the operator that appears in the LL analysis of Kovner and Lublinsky~\cite{Kovner:2006wr,Kovner:2006ge}.

\paragraph{Momentum space and $k_T$ factorisation.} With $f(\v k)=\int_xe^{i\v k\cdot\v x}f(\v x)$ and Eq.~\eqref{eq:FT_K} of Appendix~\ref{app:fourier}, the LO amplitude takes the form
\begin{align}\label{eq:ALOmom}
\mathcal A^{a}_{{\rm LO}\,i}(k)=\frac{\sqrt2\,g}{\sqrt{k^+}}\int\frac{d^2q}{(2\pi)^2}\,\Big[\frac{k^i}{\v k^2}-\frac{q^i}{\v q^2}\Big]\,U^{ab}(\v q-\v k)\,\rho^b(-\v q)\,,\qquad
U^{ab}(\v l)=\int_xe^{i\v l\cdot\v x}U^{ab}(\v x)\,,
\end{align}
so that the gluon is produced with the Lipatov-type vertex $\v k/\v k^2-\v q/\v q^2$, where $\v q$ is the momentum supplied by the projectile and $\v k-\v q$ the momentum transferred by the target. The term $U^{ab}(\v l)\propto\delta^{ab}\delta^{(2)}(\v l)$ of a trivial target does not contribute, since the vertex vanishes at $\v q=\v k$. For a dilute target, $U^{ab}(\v x)=\delta^{ab}+f^{acb}\alpha^c(\v x)+O(\alpha^2)$, the amplitude is linear in $\alpha$, and with $\big|\frac{\v k}{\v k^2}-\frac{\v q}{\v q^2}\big|^2=\frac{(\v k-\v q)^2}{\v k^2\v q^2}$ the spectrum reduces to the $k_T$-factorised form~\cite{Gribov:1983ivg,Kovchegov:2001sc,Kharzeev:2003wz,Blaizot:2004wu}. For a projectile and a target that are homogeneous in the transverse plane, with $\langle\rho^{b'}(\v q')\rho^b(-\v q)\rangle=\delta^{bb'}(2\pi)^2\delta^{(2)}(\v q-\v q')P(\v q)$ and $\langle\alpha^{c}(\v l)\alpha^{d}(-\v l')\rangle=\delta^{cd}(2\pi)^2\delta^{(2)}(\v l-\v l')T(\v l)$,
\begin{align}\label{eq:kTfact}
\frac{d^3N_{\rm LO}}{d^3k}\simeq\frac{\alpha_s N_c(N_c^2-1)}{\pi^2}\,\frac{S_\perp}{k^+}\int\frac{d^2q}{(2\pi)^2}\,\frac{(\v k-\v q)^2}{\v k^2\,\v q^2}\,P(\v q)\,T(\v k-\v q)\,,
\end{align}
where $S_\perp$ is the transverse area.

\subsection{Next-to-leading order: amplitudes}\label{sec:NLOamplitudes}

At order $g^3$ the one-gluon amplitude receives five types of contributions. Collecting the terms of $\bar\Omega^\dagger\Omega$ that connect $\ket0$ with $\bra{g}$ at this order we find
\begin{align}\label{eq:A3amp}
\mathcal A^a_{(3)\,i}(k)=\int_w e^{-i\v k\cdot\v w}\Big\{\mathcal A_{N}+\mathcal A_{A^{(3)}}+\mathcal A_{B_3}+\mathcal A_{AB_2}+\mathcal A_{CB_2}\Big\}^a_i\,,
\end{align}
with
\begin{align}
\mathcal A_{N}&=-\bar{\mf A}^a_i(k^+,-\v w)\,\mf N_2+U^{ab}(\v w)\,\bar{\mf N}_2\,\mf A^b_i(k^+,-\v w)-\big[\bar{\mf N}_2,\bar{\mf A}^a_i(k^+,-\v w)\big]\,,
\notag\\
\mathcal A_{A^{(3)}}&=U^{ab}(\v w)\,\mf A^{(3)b}_i(k^+,-\v w)-\bar{\mf A}^{(3)a}_i(k^+,-\v w)\,,
\notag\\
\mathcal A_{B_3}&=\int\frac{dp^+}{2\pi}\int_x\bar{\mf B}^{\dagger ba}_{3ji}(k^+,\v w;p^+,-\v x)\Big[U^{bc}(\v x)\,\mf A^c_j(p^+,-\v x)-\bar{\mf A}^b_j(p^+,-\v x)\Big]\,,
\notag\\
\mathcal A_{AB_2}&=2\int\frac{dp^+}{2\pi}\int_x\Big[U^{ad}(\v w)U^{bc}(\v x)\,\bar{\mf A}^{\dagger b}_j(p^+,-\v x)\,\mf B^{cd}_{2ji}(k^+,-\v w;p^+,-\v x)
\notag\\&\hspace{2.2cm}
-\bar{\mf A}^{\dagger b}_j(p^+,-\v x)\,\bar{\mf B}^{ba}_{2ji}(k^+,-\v w;p^+,-\v x)\Big]\,,
\notag\\
\mathcal A_{CB_2}&=2\int\frac{dp^+}{2\pi}\int_{x,z}\Big[\bar{\mf C}^{\dagger dca}_{1jki}(k^+,\v w-\v z;p^+,\v z-\v x)\,U^{db}(\v x)U^{ce}(\v z)\,\mf B^{be}_{2jk}(k^+-p^+,-\v z;p^+,-\v x)
\notag\\&\hspace{2.2cm}
-\bar{\mf C}^{\dagger bca}_{1jki}(k^+,\v w-\v z;p^+,\v z-\v x)\,\bar{\mf B}^{bc}_{2jk}(k^+-p^+,-\v z;p^+,-\v x)\Big]\,.
\end{align}
\paragraph{What the five terms describe.} The term $\mathcal A_N$ contains the normalisation: the term $\bar{\mf A}^{(3)}_1$ of $\bar\Omega^\dagger$ contributes, through the unitarity relation~\eqref{eq:A31unit}, the commutator $-[\bar{\mf N}_2,\bar{\mf A}]$. $\mathcal A_{A^{(3)}}$ is the $O(g^3)$ one-gluon amplitude of the wave function, before and after the collision. $\mathcal A_{B_3}$ describes the scattering of an already emitted gluon of momentum $p^+$ into the observed gluon by its interaction with the valence charges after the collision ($\bar{\mf B}_3^\dagger$); $\mathcal A_{AB_2}$, the absorption, after the collision, of one of the two gluons of the two-gluon component of the wave function; $\mathcal A_{CB_2}$, the recombination of the two gluons of that component into the observed one. In each term the two contributions correspond to an interaction of the soft gluons with the target before and after the respective valence interaction; their difference vanishes for $U=1$.

\paragraph{Nuance: the commutator term of $\mathcal A_N$ and why it cancels.} For commuting charges the term $-[\bar{\mf N}_2,\bar{\mf A}]$ of $\mathcal A_N$ would be absent. Here it is not zero: $\mf N_2$ is bilinear in the charges at two different points, and its commutator with the single charge of $\mf A$ is a two-$\rho$ term $\propto g^3f^{abc}\{\rho^b,\rho^c\}\log(\vee/\Lambda)$. Taken alone it would add a two-$\rho$ term with the full logarithm to the cross section. It does not, because the three-$\rho$ part of $\mf A^{(3)}$ has the order $\mf A\,\mf N_2$ (Section~\ref{sec:coefficients}): the terms of $\mathcal A_N$ and the three-$\rho$ part of $\mathcal A_{A^{(3)}}$ are
\begin{align}\label{eq:Nterms_steps}
&-\bar{\mf A}\,\mf N_2+U\,\bar{\mf N}_2\,\mf A-\big[\bar{\mf N}_2,\bar{\mf A}\big]+U\,\mf A\,\mf N_2-\bar{\mf A}\,\bar{\mf N}_2\,,\qquad
-\big[\bar{\mf N}_2,\bar{\mf A}\big]-\bar{\mf A}\,\bar{\mf N}_2=-\bar{\mf N}_2\,\bar{\mf A}\,,
\end{align}
so that the commutator only reverses the order of one product. With the ordering of the three-$\rho$ part of $\mf A^{(3)}$ derived in Section~\ref{sec:coefficients}, $[\mf A^{(3)}]_{\rho\rho\rho}=\mf A\,\mf N_2$, the commutator in $\mathcal A_N$ cancels against the corresponding part of $\mathcal A_{A^{(3)}}$, and the four-$\rho$ (normalisation) terms of the amplitude add up to
\begin{align}\label{eq:Nterms}
-\bar{\mf A}\,\mf N_2+U\,\bar{\mf N}_2\,\mf A+U\,\mf A\,\mf N_2-\bar{\mf N}_2\,\bar{\mf A}\,,
\end{align}
which vanishes for $U=1$, as it must, and in which no commutator is left. With the opposite ordering, $\mf N_2\mf A$, a spurious two-$\rho$ term $\propto\log(\vee/\Lambda)$ would remain, which does not vanish for $U\neq1$. The explicit form of the commutator shows what is at stake,
\begin{align}\label{eq:commN2A}
\big[\mf N_2,\mf A^a_i(k^+,-\v w)\big]=-\frac{g^3}{8\sqrt2\,\pi^4\sqrt{k^+}}\,\log\frac{\vee}{\Lambda}\,f^{abc}\int_{x,y,z}K(\v x,\v y,\v z)\,\Kv^i(\v y-\v w)\,\{\rho^b(\v x),\rho^c(\v y)\}\,:
\end{align}
it is equal, with the same coefficient and sign, to the term $-2k^i\log(\vee/\Lambda)$ of the two-$\rho$ part of $\mf A^{(3)}$, Eq.~\eqref{eq:A3tworho} (we verified this also with explicit $SU(2)$ and $SU(3)$ matrices). The two descriptions are therefore equivalent: either the three-$\rho$ term is ordered as $\mf A\,\mf N_2$ and the two-$\rho$ part~\eqref{eq:A3tworho} is used as it stands, which is what we do, or the three-$\rho$ term is written in the opposite order and the term $-2k^i\log(\vee/\Lambda)$ is removed from the unbarred two-$\rho$ structure. What is not allowed is to combine the opposite order with Eq.~\eqref{eq:A3tworho}.

The two-gluon amplitude at order $g^2$ is
\begin{align}\label{eq:A2gamp}
\mathcal A^{ab}_{ij}(k;q^+,\v v)=\frac{2}{\sqrt{2\pi}}\int_w e^{-i\v k\cdot\v w}\Big\{\mathcal A_B+\mathcal A^{(1)}_{AA}+\mathcal A^{(2)}_{AA}-\mathcal A_{CA}\Big\}^{ab}_{ij}\,,
\end{align}
with
\begin{align}
\mathcal A_B&=U^{ac}(\v w)\,\mf B^{bc}_{2ji}(k^+,-\v w;q^+,-\v v)-U^{\dagger bc}(\v v)\,\bar{\mf B}^{ca}_{2ji}(k^+,-\v w;q^+,-\v v)\,,
\notag\\
\mathcal A^{(1)}_{AA}&=\tfrac12\Big[U^{\dagger bc}(\v v)\,\bar{\mf A}^c_j(q^+,-\v v)\,\bar{\mf A}^a_i(k^+,-\v w)-U^{ac}(\v w)\,\bar{\mf A}^d_j(q^+,-\v v)\,U^{\dagger bd}(\v v)\,\mf A^c_i(k^+,-\v w)\Big]\,,
\notag\\
\mathcal A^{(2)}_{AA}&=\tfrac12\Big[U^{\dagger bc}(\v v)\,\bar{\mf A}^a_i(k^+,-\v w)\,\bar{\mf A}^c_j(q^+,-\v v)-\bar{\mf A}^a_i(k^+,-\v w)\,\mf A^b_j(q^+,-\v v)\Big]\,,
\notag\\
\mathcal A_{CA}&=\int_z\Big[\bar{\mf C}^{dac}_{1jik}(q^++k^+,\v z-\v w;q^+,\v w-\v v)\,U^{\dagger bd}(\v v)\,U^{ce}(\v z)\,\mf A^e_k(q^++k^+,-\v z)
\notag\\&\hspace{1.2cm}
-U^{\dagger bd}(\v v)\,\bar{\mf C}^{dac}_{1jik}(q^++k^+,\v z-\v w;q^+,\v w-\v v)\,\bar{\mf A}^c_k(q^++k^+,-\v z)\Big]\,.
\end{align}
$\mathcal A_B$ is the two-gluon component of the wave function scattered on the target; $\mathcal A^{(1)}_{AA}$ and $\mathcal A^{(2)}_{AA}$ arise from the product of two single emissions, one before and one after the collision, in the two operator orderings generated by $\bar\Omega^\dagger$; and $\mathcal A_{CA}$ is the splitting, after the collision, of a gluon of momentum $q^++k^+$ emitted before the collision into the observed gluon and a second gluon of momentum $q^+$.

\subsection{Next-to-leading order: Groups and rows}\label{sec:groups}

\paragraph{Assembly.} Inserting $\mathcal A=\mathcal A_{\rm LO}+\mathcal A_{(3)}$ and the two-gluon amplitude~\eqref{eq:A2gamp} into Eq.~\eqref{eq:N_amplitudes} and keeping the terms of order $g^4$, the one-gluon final state gives $\mathcal A_{\rm LO}^\dagger\mathcal A_{(3)}+{\rm c.c.}$ (the square of $\mathcal A_{(3)}$ is of order $g^6$), and the two-gluon final state gives the square of the sum of the four parts of $\mathcal A^{ab}_{ij}$, i.e.\ four squares and six interferences. The $O(g^4)$ cross section is
\begin{align}\label{eq:NLOdecomp}
\frac{d^3N}{d^3k}=\frac{d^3N_{\rm LO}}{d^3k}+\frac{d^3N^{\rm(I)}}{d^3k}+\frac{d^3N^{\rm(II)}}{d^3k}+\frac{d^3N^{\rm(III)}}{d^3k}+O(g^6)\,,
\end{align}
where
\begin{align}\label{eq:Groups}
\frac{d^3N^{\rm(I)}}{d^3k}&=\frac{1}{(2\pi)^3}\Big\langle\mathcal A^{a\dagger}_{{\rm LO}\,i}(k)\,\mathcal A^a_{(3)\,i}(k)\Big\rangle+{\rm c.c.}\,,\notag\\
\frac{d^3N^{\rm(II)}}{d^3k}&=\frac{1}{(2\pi)^3}\int dq^+\!\int_v\Big\langle\big|\mathcal A_B\big|^2+\big|\mathcal A^{(1)}_{AA}\big|^2+\big|\mathcal A^{(2)}_{AA}\big|^2+\big|\mathcal A_{CA}\big|^2\Big\rangle\,,\notag\\
\frac{d^3N^{\rm(III)}}{d^3k}&=\frac{1}{(2\pi)^3}\int dq^+\!\int_v\Big\langle\mathcal A^{(1)\dagger}_{AA}\mathcal A_B+\mathcal A^{(2)\dagger}_{AA}\mathcal A_B-\mathcal A^{\dagger}_{CA}\mathcal A_B
\notag\\&\hspace{3.3cm}+\mathcal A^{(2)\dagger}_{AA}\mathcal A^{(1)}_{AA}-\mathcal A^{\dagger}_{CA}\mathcal A^{(1)}_{AA}-\mathcal A^{\dagger}_{CA}\mathcal A^{(2)}_{AA}\Big\rangle+{\rm c.c.}
\end{align}
Here $|\mathcal A|^2\equiv\mathcal A^{ab\dagger}_{ij}\mathcal A^{ab}_{ij}$, the sign of $\mathcal A_{CA}$ is that of Eq.~\eqref{eq:A2gamp}, and the momentum of the integrated gluon is renamed $q^+\to p^+$ below. Group~I is the interference of the LO amplitude with the $O(g^3)$ one-gluon amplitude; Group~II consists of the squares of the four components of the two-gluon amplitude; Group~III of their interferences. Groups~I and~III are written with their complex conjugates; Group~II is Hermitian by itself. We call the individual products of coefficients \emph{rows}. They are listed in Table~\ref{tab:rows}, together with the range of the longitudinal integration, the number of charges and the parts of the final result to which they contribute.

\begin{table}[t]
\centering
\renewcommand{\arraystretch}{1.2}
\footnotesize
\begin{tabular}{@{}l p{4.0cm} p{2.9cm} l p{3.4cm}@{}}
\toprule
Row & product (Eqs.~\eqref{eq:A3amp},~\eqref{eq:A2gamp}) & range of $p^+$ & $\#\rho$ & contributes to\\
\midrule
I.I & LO$^\dagger\times$ normalisation terms~\eqref{eq:Nterms} & --- & 4 & 4$\rho$ ($\log\frac{\vee}{\Lambda}$)\\
I.II & LO$^\dagger\times[U\mf A^{(3)}-\bar{\mf A}^{(3)}]$ & --- & 2--4 & 4$\rho$, BFKL, JIMWLK, 3$\rho$, finite\\
I.III & LO$^\dagger\times\bar{\mf B}_3^\dagger\mf A$ & $\Lambda<p^+<\vee$, $|p^+-k^+|>\Lambda$ & 3, 4 & BFKL, 3$\rho$, finite\\
I.IV & LO$^\dagger\times\bar{\mf A}^\dagger\mf B_2$ & $\Lambda<p^+<\vee-k^+$ $[\vee]$ & 3, 4 & 4$\rho$, BFKL, 3$\rho$; finite part $=-$III.I\\
I.V & LO$^\dagger\times\bar{\mf C}_1^\dagger\mf B_2$ & $\Lambda<p^+<k^+-\Lambda$ & 2, 3 & JIMWLK, 3$\rho$, finite\\
\midrule
II.I & $|\mathcal A_B|^2$ & $\Lambda<p^+<\vee-k^+$ $[\vee]$ & 2--4 & 4$\rho$, BFKL, JIMWLK, 3$\rho$, finite\\
II.II & $|\mathcal A^{(1)}_{AA}|^2$ & $\Lambda<p^+<\vee$ & 4 & 4$\rho$\\
II.III & $|\mathcal A^{(2)}_{AA}|^2$ & $\Lambda<p^+<\vee$ & 4 & 4$\rho$\\
II.IV & $|\mathcal A_{CA}|^2$ & $\Lambda<p^+<\vee-k^+$ & 2 & BFKL, JIMWLK, DGLAP, RC, finite\\
\midrule
III.I & $\mathcal A^{(1)\dagger}_{AA}\mathcal A_B$ & $\Lambda<p^+<\vee-k^+$ $[\vee]$ & 3, 4 & 4$\rho$, BFKL, 3$\rho$; finite part $=-$I.IV\\
III.II & $\mathcal A^{(2)\dagger}_{AA}\mathcal A_B$ & $\Lambda<p^+<\vee-k^+$ $[\vee]$ & 3, 4 & 4$\rho$, BFKL, 3$\rho$, finite\\
III.III & $-\mathcal A^{\dagger}_{CA}\mathcal A_B$ & $\Lambda<p^+<\vee-k^+$ & 2, 3 & BFKL, JIMWLK, 3$\rho$, finite\\
III.IV & $\mathcal A^{(2)\dagger}_{AA}\mathcal A^{(1)}_{AA}$ & $\Lambda<p^+<\vee$ & 4 & 4$\rho$\\
III.V & $-\mathcal A^{\dagger}_{CA}\mathcal A^{(1)}_{AA}$ & $\Lambda<p^+<\vee-k^+$ & 3 & BFKL, 3$\rho$, finite\\
III.VI & $-\mathcal A^{\dagger}_{CA}\mathcal A^{(2)}_{AA}$ & $\Lambda<p^+<\vee-k^+$ & 3 & BFKL, 3$\rho$, finite\\
\bottomrule
\end{tabular}
\caption{The rows of the $O(g^4)$ cross section. ``LO$^\dagger$'' is the conjugate LO amplitude, Eq.~\eqref{eq:ALO}; $p^+$ is the longitudinal momentum of the second (integrated or virtual) gluon. Where two ranges apply, the first refers to the part linear in $\rho$ of $\mf B_2$ (its parent gluon carries $p^++k^+$) and the bracket $[\vee]$ to its $\{\rho,\rho\}$ part, $\Lambda<p^+<\vee$. ``$\#\rho$'' is the number of charges before reordering; ``4$\rho$'' and ``3$\rho$'' denote the terms quartic and cubic in $\rho$ multiplying a rapidity logarithm, which cancel (Section~\ref{sec:rapidity}). Row I.II contains in addition the ultraviolet $\beta_0$ terms of the vertex correction, which cancel against Row I.V (Section~\ref{sec:dglap_rc}).}
\label{tab:rows}
\end{table}

A glossary of the rows, with the physical process that each of them describes, is given in Table~\ref{tab:glossary}. The diagrams of the wave-function components from which the rows are built are those of Appendix~\ref{app:wavefunctions}; the last column of the table lists them. We do not draw separate diagrams for the rows: a row is the product of two such components, one of them conjugated, with the shock wave of the target inserted before or after each valence interaction.

\begingroup
\footnotesize
\renewcommand{\arraystretch}{1.25}
\begin{longtable}{@{}l>{\raggedright\arraybackslash}p{0.11\textwidth}>{\raggedright\arraybackslash}p{0.56\textwidth}>{\raggedright\arraybackslash}p{0.15\textwidth}@{}}
\caption{Glossary of the rows: type, physical content, and the figures with the diagrams of the wave-function components involved. ``Virtual'': one gluon in the final state, the second gluon is emitted and reabsorbed; ``real'': two gluons in the final state, the second is integrated over. In all rows the measured gluon is at $\v w$ in the amplitude and at $\v w'$ in the conjugate amplitude.}\label{tab:glossary}\\
\toprule
Row & type & what interferes, and which gluon is measured & components\\
\midrule
\endfirsthead
\toprule
Row & type & what interferes, and which gluon is measured & components\\
\midrule
\endhead
\bottomrule
\endfoot
I.I & virtual, $4\rho$ & LO$^\dagger$ with the LO emission multiplied by the probability amplitude $\mf N_2$ that no second soft gluon is present, before or after the collision. Measured: the LO gluon. & Fig.~\ref{fig:vac_g}\\
I.II & virtual & LO$^\dagger$ with the $O(g^3)$ one-gluon component of the wave function, scattered on the target or subtracted after it: self-energy of the measured gluon and its vertex corrections on the valence charges. Measured: the gluon of that component. & Figs.~\ref{fig:vac_g3_onerho}--\ref{fig:vac_g3_threerho}\\
I.III & virtual & LO$^\dagger$ with a LO gluon of momentum $p^+$ that, after the collision, is turned into the measured gluon of momentum $k^+$ by an interaction with the rotated valence charges ($\bar{\mf B}_3^\dagger$): final-state rescattering on the projectile. Measured: the gluon after this transition. & Figs.~\ref{fig:vac_g}, \ref{fig:g_g}\\
I.IV & virtual & LO$^\dagger$ with the two-gluon component, one of whose gluons is absorbed by the valence charges after the collision ($\bar{\mf A}^\dagger$). Measured: the other gluon. & Figs.~\ref{fig:vac_g}, \ref{fig:vac_gg}\\
I.V & virtual & LO$^\dagger$ with the two-gluon component whose gluons recombine, after the collision, into one ($\bar{\mf C}_1^\dagger$). Measured: the recombined gluon, at the $p^+$-weighted centre of the pair. & Figs.~\ref{fig:vac_gg}, \ref{fig:g_gg}\\
\midrule
II.I & real & Square of the two-gluon component of the wave function scattered on the target. Measured: one of the two gluons. & Fig.~\ref{fig:vac_gg}\\
II.II & real, $4\rho$ & Square of two independent emissions, one before and one after the collision, first operator ordering. & Fig.~\ref{fig:vac_g}\\
II.III & real, $4\rho$ & The same in the second operator ordering. & Fig.~\ref{fig:vac_g}\\
II.IV & real & Square of the amplitude in which a LO gluon of momentum $p^++k^+$ splits after the collision ($\bar{\mf C}_1$). Measured: one daughter; its parent sits at $\v y=\v z+\xi(\v w-\v z)$. The row with the collinear divergence. & Figs.~\ref{fig:vac_g}, \ref{fig:g_gg}\\
\midrule
III.I & real & Independent emissions (first ordering)$^\dagger$ $\times$ two-gluon component. & Figs.~\ref{fig:vac_g}, \ref{fig:vac_gg}\\
III.II & real & Independent emissions (second ordering)$^\dagger$ $\times$ two-gluon component. & Figs.~\ref{fig:vac_g}, \ref{fig:vac_gg}\\
III.III & real & Final-state splitting$^\dagger$ $\times$ two-gluon component: the pair is produced by splitting after the collision in the conjugate amplitude and is already present in the wave function in the amplitude. & Figs.~\ref{fig:vac_gg}, \ref{fig:g_gg}\\
III.IV & real, $4\rho$ & Interference of the two orderings of the independent emissions. & Fig.~\ref{fig:vac_g}\\
III.V & real & Final-state splitting$^\dagger$ $\times$ independent emissions (first ordering). Measured: a daughter in the conjugate amplitude (recoil phase with $\v y'$), an independently emitted gluon in the amplitude. & Figs.~\ref{fig:vac_g}, \ref{fig:g_gg}\\
III.VI & real & The same with the second ordering of the independent emissions. & Figs.~\ref{fig:vac_g}, \ref{fig:g_gg}\\
\end{longtable}
\endgroup

\paragraph{Nuance: kinematics, not the form of a coefficient, fixes the limits of the $p^+$ integrals.} The longitudinal ranges in Table~\ref{tab:rows} follow from the kinematics of the coefficients. All soft gluons satisfy $\Lambda<p^+<\vee$. Whenever a soft gluon of momentum $p^++k^+$ appears as an intermediate state (the parent of $\mf C_1$, the part of $\mf B_2$ linear in $\rho$), the condition $p^++k^+<\vee$ applies. The $\{\rho,\rho\}$ part of $\mf B_2$ describes two independent emissions and has support on the full square $\Lambda<k^+,p^+<\vee$; this is required for the cancellation of the four-$\rho$ terms (Section~\ref{sec:fourrho}). In $\mf B_3$ the momentum difference of the incoming and outgoing gluon is carried by the valence charges and satisfies $|p^+-k^+|>\Lambda$. These supports have to be read off from the kinematics of each diagram and cannot be copied from the form in which a coefficient happens to be written. A support such as $p^+<k^+-\Lambda$ for $\mf B_2$, for instance, would not be symmetric under the exchange of its two gluons and would remove the region $p^+>k^+$, which is the region that produces the projectile logarithm $\log(\vee/k^+)$; and the two supports of the two parts of $\mf B_2$ decide whether the four-$\rho$ terms cancel. Since the range of $p^+$ decides which logarithm a row produces (Section~\ref{sec:loglong}), an error in a support is an error in the evolution.

The evaluation of each row consists of the colour algebra of the products of coefficients, with the charges kept in their order, the longitudinal integration over $p^+$, and the transverse integrations that can be done in closed form. Three kinds of terms emerge:
\begin{enumerate}
\item terms proportional to $\log(\vee/k^+)$ or $\log(k^+/\Lambda)$, quadratic, cubic or quartic in $\rho$, which arise from the regions $p^+\gg k^+$ and $p^+\ll k^+$ of the longitudinal integrals. They are discussed in Section~\ref{sec:rapidity};
\item terms with a collinear singularity, from the region in which the integrated gluon of Row~II.IV is collinear with the observed one, and terms proportional to $\beta_0$ (Section~\ref{sec:dglap_rc});
\item the remaining finite terms (Section~\ref{sec:finite}).
\end{enumerate}
Two rows require a remark. The non-logarithmic parts of Rows~I.IV and III.I cancel identically. Rows~II.II, II.III and III.IV are products of four coefficients $\mf A$, each linear in $\rho$; they contain only four-$\rho$ terms.

\begin{summarybox}
\begin{itemize}
\item The observable is the number of dressed gluons in the scattered state, Eq.~\eqref{eq:observable}; it is a sum of squared one- and two-gluon amplitudes of the operator $\bar\Omega^\dagger\Omega$, Eq.~\eqref{eq:N_amplitudes}, which is the identity without a target.
\item LO: Eq.~\eqref{eq:LO}; for a colour-diagonal projectile Eq.~\eqref{eq:LOsinglet}, the formula of Kovchegov and Tuchin. The colour structure is a product of two differences $U(\text{gluon})-U(\text{source})$.
\item NLO: the $O(g^3)$ one-gluon amplitude has five parts, Eq.~\eqref{eq:A3amp}, the $O(g^2)$ two-gluon amplitude four, Eq.~\eqref{eq:A2gamp}. Their products are the fifteen rows of Tables~\ref{tab:rows} and~\ref{tab:glossary}: Group~I (virtual), Groups~II and~III (real).
\item The commutator term required by unitarity cancels against the ordered three-$\rho$ term of $\mf A^{(3)}$; the range of each $p^+$ integral is fixed by kinematics.
\end{itemize}
\end{summarybox}
